\section{Related Work}

\subsection{Information-Theoretic Causality: Transfer Entropy}
The inference of directed causal relationships from observational time series traces back to Granger's predictive causality \cite{granger1969investigating} and Pearl's structural causal diagrams \cite{pearl2009causality}. While Granger causality relies on linear autoregressive modeling, complex physiological and biological systems exhibit prominent non-linear couplings \cite{runge2019inferring}. To overcome this limitation, Schreiber formulated Transfer Entropy (TE) \cite{schreiber2000measuring}, a model-free metric that quantifies the asymmetric exchange of information by measuring the conditional mutual information between the source history and the target's current state.

Transfer Entropy has been widely applied to discover effective connectivity in neuroimaging \cite{vicente2011transfer} and to quantify information transfer delays in autonomic regulation \cite{wibral2013measuring}. Traditionally, TE is estimated using non-parametric methods, with the Kraskov-St{\"o}gbauer-Grassberger (KSG) estimator \cite{kraskov2004estimating} serving as the standard reference. By constructing adaptive rectangular neighborhoods via $k$-nearest neighbor ($k$-NN) statistics, KSG cancels out marginal bias in the asymptotic limit ($N \to \infty$). However, physiological monitoring often restricts analysis to short, locally stationary observation windows ($N \le 100$). In this sparse regime, $k$-NN distance distributions become highly erratic due to distance concentration in higher dimensions ($d \ge 3$), inducing severe estimation variance that obscures genuine physiological coupling.

\subsection{Neural Estimators and the Modern SOTA Landscape}
To overcome the limitations of distance-based density estimation, deep learning has been increasingly incorporated into information-theoretic analysis \cite{fawaz2019deep}. A milestone in this direction was Mutual Information Neural Estimation (MINE) \cite{belghazi2018mine}, which leverages the Donsker-Varadhan (DV) variational characterization of Kullback-Leibler (KL) divergence \cite{donsker1975asymptotic} to train a parameterised statistics network via backpropagation.

Recently, two distinct state-of-the-art frameworks have been introduced for time-series transfer entropy:
\begin{enumerate}
    \item \textbf{TREET (Transformer-based TE, 2024/2025) \cite{treet2024}:} TREET parameterizes the statistics network using multi-head causal attention and derives theoretical convergence through Birkhoff's ergodic theorem. While effective at extracting long-memory contexts ($l = 99$ samples), TREET relies on dense sequence interactions. On ultra-short physiological windows ($N \le 50$), attention matrices undergo severe dispersion, leading to negative transfer entropy estimates (e.g., $-3.02 \pm 2.59$ on uncoupled synthetic benchmarks) and high variance.
    \item \textbf{TENDE (Score-based Diffusion SDE, 2026) \cite{tende2025}:} TENDE addresses the fundamental theoretical limitation identified by McAllester and Stratos \cite{mcallester2020formal}—which proves that variational lower bounds require an exponential sample size $\mathcal{O}(e^I)$ to converge for large information values $I$. By establishing a closed-form connection between KL divergence and diffusion score matching under Variance-Preserving Stochastic Differential Equations (VP-SDE), TENDE sidesteps variational lower bounds entirely. However, numerically integrating the reverse-time SDE across hundreds of denoising steps requires approximately 20 minutes per time-series pair on NVIDIA A100 GPUs, rendering it completely unfeasible for continuous wearable or edge deployment.
\end{enumerate}

\subsection{Amortized Inference and the Reality Gap in Biosignals}
An efficient alternative to per-sample optimization is \emph{amortized inference}, wherein a neural network is pre-trained across a large synthetic distribution corpus and executes inferences via a single forward pass ($\mathcal{O}(1)$ complexity). However, transferring amortized neural estimators to real-world biomedical signals encounters an acute \emph{reality gap} (covariate shift) \cite{iyengar1996age}.

While classical Unsupervised Domain Adaptation (UDA) frameworks in computer vision utilize gradient reversal layers \cite{ganin2015unsupervised} or adversarial domain discriminators \cite{tzeng2017adversarial}, these adversarial objectives are unstable and prone to mode collapse when applied to sparse time series. In this work, we bypass complex adversarial architectures by exploiting the variational Donsker-Varadhan objective directly at test time. Because the DV bound operates on temporally permuted marginal distributions, it allows the network to adapt to patient-specific manifolds in a purely unsupervised, label-free manner.

\subsection{Quantum Machine Learning for Dynamical Interactions}
To push beyond classical neural representational bounds, Variational Quantum Circuits (VQC) parameterized on near-term quantum processors (NISQ) have emerged as an expressive class of machine learning models \cite{bergholm2018pennylane}. Through data re-uploading techniques \cite{perez2020data}, low-dimensional classical coordinates are mapped into high-dimensional complex Hilbert spaces $\mathbb{C}^{2^n}$ where non-linear unitary transformations and entangled states can capture intricate dependency structures. In this paper, we explore the integration of VQC within the amortized TE framework, providing both theoretical and empirical insights into hybrid classical-quantum causal estimation.
